\chapter*{Preface}
\addcontentsline{toc}{chapter}{Preface}

This book is an experiment.

Category theory has a reputation: it is something you learn \emph{after} years
of mathematics, as a kind of universal solvent that dissolves the walls between
algebra, geometry, logic, and the rest. Every existing textbook on the subject
assumes you already live in those rooms.

This book tries something different. It treats category theory not as a
retrospective unifier of mathematics but as a \emph{starting point}. We build
from nothing. The only raw material is the idea that there are \term{things},
and between things there are \term{connections}. Everything else grows from
there.

\bigskip

\textbf{Who this book is for.} If you can follow a recipe, read a subway map,
or understand why two different roads can take you to the same place, you have
all the prerequisites. We will not assume familiarity with algebra, calculus,
set theory, or any other branch of mathematics. When those subjects appear
later as examples, they will be introduced fresh.

\bigskip

\textbf{How this book is organized.} Each section is written in four layers,
labeled clearly:

\begin{itemize}
  \item \colorbox{CoreColor}{\color{white}\footnotesize\bfseries\ CORE\ } ---
    The essential idea in plain language. Read this even if you read nothing
    else in the section.

  \item \colorbox{Exp1Color}{\color{white}\footnotesize\bfseries\ EXPANSION 1\ } ---
    Worked examples, diagrams, and the kind of out-loud thinking you should
    practice when you encounter these ideas in the wild.

  \item \colorbox{Exp2Color}{\color{white}\footnotesize\bfseries\ EXPANSION 2\ } ---
    A deeper look: variations, subtleties, and the ``what if'' questions that
    reward careful readers.

  \item \colorbox{Exp3Color}{\color{white}\footnotesize\bfseries\ EXPANSION 3\ } ---
    Formal treatment. Precise definitions, proofs, and the vocabulary that
    connects this material to the broader literature.
\end{itemize}

You can read straight through, or you can read only the Core sections of every
chapter first and then go back for depth. Both paths work. The book will not
punish you for skipping an Expansion layer the first time through.

\bigskip

\textbf{The internal-dialog boxes.} Throughout the book you will find boxes
labeled ``Internal Dialog.'' These model the running commentary that a skilled
reader maintains while working with diagrams and definitions. Learning to
produce that commentary yourself --- in your own head, in pencil in the margin
--- is one of the main skills this book tries to teach. Mathematical notation
is dense; the internal dialog is what keeps it from being opaque.

\bigskip

\textbf{A note on what we avoid.} We will not say things like ``consider the
category of groups'' or ``the fundamental groupoid of a topological space''
until we have built those ideas from scratch. The temptation to borrow
ready-made examples from other fields is real, but it defeats the purpose. Our
examples will come from everyday experience: maps and routes, instructions and
steps, classifications and comparisons. The mathematics, when it appears, will
be mathematics we have built together.

\bigskip

Let's begin.
